\section{Introduction}
\IEEEPARstart{L}{arge} language model (LLM) agents increasingly participate in
industrial operation: they generate operator-reviewed OPC UA control sequences
\cite{huang2026made}, invoke factory skills over tool protocols
\cite{koehle2026llmdriven}, coordinate shop-floor tasks
\cite{lim2025llmenabled,zhao2026shopfloor}, and schedule a simulated
robot-operated line \cite{waseem2026pretrained}. These applications require
a connection between reasoning and production context: who may change a
parameter, under which product and recipe, what was commanded, and how the
process responded.

To bring injection-molded part weight into specification, an agent may use
recent measurements to diagnose insufficient packing and propose higher
holding pressure. Execution also requires a writable signal mapping,
parameter limits, authorized production scope, and approval. Its record
must distinguish the pressure request, endpoint confirmation, and measured
weight.

Three distinctions must survive this handoff: permission to invoke a tool
versus authority over a production
object, transport acceptance versus endpoint confirmation, and a trial
command versus a retained recipe. Reconstructing these distinctions inside
every agent application couples execution policy to the selected reasoning
engine and obscures how an intervention relates to a running batch. We
address the following systems question: how can an industrial
agent platform retain a common execution and evidence contract when the
proposal strategy, equipment protocol, and production context change?

We present a governed supervisory execution framework connecting agent
teams to semantic data-acquisition (DAQ) and write-control (DCW) nodes.
A deterministic execution boundary sits beneath the replaceable strategy
layer. The strategy acquires evidence, proposes a change, and interprets
observations. The platform resolves authority, applies source-specific
limits and required approval, dispatches commands, and grades confirmation.
Line--product--recipe--run objects supply production context.
Separate task, intervention, and batch records then connect the original
intent to commanded controls and observed outputs.

The contribution is a common execution-and-evidence contract integrating
production-scoped authority, graded confirmation, and route-specific records:
\begin{enumerate}
\item a capability-based architecture that binds channel identity and
semantic industrial nodes to line--product--recipe--run context, separating
information exchange from execution authority;
\item route-specific execution mechanisms that join admission, approval,
graded confirmation, observation, and compensation while preserving
distinct node-intervention and recipe-batch records; and
\item an implementation and evaluation of this contract across replaceable
strategies and protocol mappings, using admission ablation, confirmation
drills, process missions, repeated regression runs, and a recipe trajectory.
\end{enumerate}

Sections~\ref{sec:arch}--\ref{sec:mech} explain the architecture and execution
mechanisms; Sections~\ref{sec:eval}--\ref{sec:limits} evaluate and interpret
their simulator-backed evidence.
